% --- CORRECCION: Forzar interpretacion raw de la entrada ---
\UseRawInputEncoding
\documentclass[aspectratio=169, 10pt, xcolor=table]{beamer}

% --- Configuracion del Tema ---
\usetheme{Madrid}
\usecolortheme{dolphin}

% --- Paquetes necesarios ---
\usepackage[utf8]{inputenc}
\usepackage[spanish]{babel}
\usepackage{amsmath, amssymb}
\usepackage{booktabs}
\usepackage{graphicx}
\usepackage{listings}
\usepackage{tikz}
\usetikzlibrary{arrows.meta, positioning, shapes.geometric}
\usepackage{etoolbox}

% --- Ruta de Imagenes ---
\graphicspath{{./imagenes/}}

% --- Estilo para codigo Python ---
\definecolor{codegreen}{rgb}{0,0.6,0}
\definecolor{codegray}{rgb}{0.5,0.5,0.5}
\definecolor{codepurple}{rgb}{0.58,0,0.82}
\definecolor{backcolour}{rgb}{0.95,0.95,0.92}
\lstset{
    language=Python,
    backgroundcolor=\color{backcolour},
    commentstyle=\color{codegreen},
    keywordstyle=\color{blue},
    numberstyle=\tiny\color{codegray},
    stringstyle=\color{codepurple},
    basicstyle=\ttfamily\scriptsize,
    breaklines=true,
    frame=single,
    showstringspaces=false,
    literate={á}{{\'a}}1 {é}{{\'e}}1 {í}{{\'i}}1 {ó}{{\'o}}1 {ú}{{\'u}}1 {ñ}{{\~n}}1 {ü}{{\"u}}1
}

% --- Pie de pagina personalizado ---
\makeatletter
\setbeamertemplate{footline}{
  \leavevmode%
  \hbox{%
  \begin{beamercolorbox}[wd=.333333\paperwidth,ht=2.25ex,dp=1ex,center]{author in head/foot}%
    \usebeamerfont{author in head/foot}\insertshortauthor
  \end{beamercolorbox}%
  \begin{beamercolorbox}[wd=.333333\paperwidth,ht=2.25ex,dp=1ex,center]{title in head/foot}%
    \usebeamerfont{title in head/foot}Sesion 8
  \end{beamercolorbox}%
  \begin{beamercolorbox}[wd=.333333\paperwidth,ht=2.25ex,dp=1ex,right]{date in head/foot}%
    \usebeamerfont{date in head/foot}CALCULO Y ALGEBRA LINEAL \hfill \insertframenumber/\inserttotalframenumber \hspace*{0.3cm}
  \end{beamercolorbox}}%
  \vskip0pt%
}
\makeatother

% --- Datos de la portada ---
\title[SESION 8 CALCULO]{\textbf{Sesion 8}}
\subtitle{Derivadas Parciales y Vector Gradiente}
\author[Prof. C. Sanchez]{Carlos Cesar Sanchez Coronel}
\institute{\textbf{CALCULO Y ALGEBRA LINEAL PARA IA}}
\date{}

% --- Eliminar barra de navegacion ---
\setbeamertemplate{navigation symbols}{}

% --- Indice al inicio de cada seccion ---
\AtBeginSection[]{
  \begin{frame}
    \frametitle{Contenido}
    \tableofcontents[currentsection]
  \end{frame}
}

\begin{document}

% --- Portada ---
\begin{frame}
    \titlepage
\end{frame}

% --- Indice general ---
\begin{frame}{Contenido}
    \tableofcontents
\end{frame}

% ============================================================================
% SECCION 1: INTRODUCCION A DERIVADAS PARCIALES
% ============================================================================
\section{Derivadas Parciales: Sensibilidad Individual}

\begin{frame}{Motivacion}
    \begin{itemize}
        \item En IA, las funciones dependen de muchas variables (pesos, caracteristicas).
        \item Necesitamos medir el cambio al variar una variable, manteniendo las demas constantes.
        \item De ahi nacen las \textbf{derivadas parciales}.
    \end{itemize}
\end{frame}

\begin{frame}{Definicion intuitiva}
    \begin{block}{Derivada parcial respecto a $x_i$}
        \[
        \frac{\partial f}{\partial x_i}(\mathbf{a}) = \lim_{h\to 0} \frac{f(a_1,\dots,a_i+h,\dots,a_n) - f(a_1,\dots,a_n)}{h}
        \]
    \end{block}
    \begin{itemize}
        \item Mide la tasa de cambio de $f$ cuando solo $x_i$ varia.
        \item Es la derivada ordinaria de la funcion $g(t)=f(\dots,t,\dots)$ en $t=a_i$.
    \end{itemize}
\end{frame}

\begin{frame}{Ejemplo basico}
    Sea $f(x,y)=x^2y + \sin(xy)$. Calculamos:
    \[
    \frac{\partial f}{\partial x} = 2xy + y\cos(xy),\qquad 
    \frac{\partial f}{\partial y} = x^2 + x\cos(xy)
    \]
    \vspace{0.3cm}
    \begin{itemize}
        \item Notacion: tambien $f_x$, $f_y$, $\partial_x f$, etc.
    \end{itemize}
\end{frame}

\begin{frame}{Interpretacion geometrica (2D)}
    \begin{itemize}
        \item $z=f(x,y)$ es una superficie.
        \item $\frac{\partial f}{\partial x}(a,b)$ es la pendiente de la curva al cortar con el plano $y=b$.
    \end{itemize}
    \centering
    \begin{tikzpicture}[scale=0.8]
        \draw[->] (-1,0) -- (4,0) node[right] {$x$};
        \draw[->] (0,-1) -- (0,3) node[above] {$z$};
        \draw[thick, blue] plot[smooth] coordinates {(0,0.5) (1,1) (2,1.8) (3,2.5)} node[right] {$f(x,b)$};
        \draw[dashed] (1.5,0) -- (1.5,1.4);
        \draw[dashed] (0,1.4) -- (1.5,1.4);
        \fill (1.5,1.4) circle (2pt) node[above] {$(a,f(a,b))$};
    \end{tikzpicture}
\end{frame}

% ============================================================================
% SECCION 2: VECTOR GRADIENTE
% ============================================================================
\section{Vector Gradiente: Direccion de Maximo Crecimiento}

\begin{frame}{Definicion}
    \begin{block}{Gradiente de $f:\mathbb{R}^n\to\mathbb{R}$}
        \[
        \nabla f(\mathbf{x}) = 
        \begin{bmatrix}
            \frac{\partial f}{\partial x_1}(\mathbf{x}) \\[4pt]
            \frac{\partial f}{\partial x_2}(\mathbf{x}) \\[4pt]
            \vdots \\[4pt]
            \frac{\partial f}{\partial x_n}(\mathbf{x})
        \end{bmatrix}
        \]
    \end{block}
\end{frame}

\begin{frame}{Interpretacion geometrica}
    \begin{itemize}
        \item Apunta en la direccion de \textbf{maximo crecimiento} de la funcion.
        \item Su magnitud $\|\nabla f\|$ indica la tasa de crecimiento en esa direccion.
        \item Es perpendicular a las curvas de nivel (lineas de $f$ constante).
    \end{itemize}
    \centering
    \includegraphics[width=0.5\textwidth]{gradiente_curvas_nivel.png} % si existe
\end{frame}

\begin{frame}{Propiedades importantes}
    \begin{itemize}
        \item Linealidad: $\nabla(\alpha f + \beta g) = \alpha\nabla f + \beta\nabla g$.
        \item Regla del producto: $\nabla(fg) = f\nabla g + g\nabla f$.
        \item Regla de la cadena multivariable: 
        \[
        \frac{d}{dt} f(g(t)) = \nabla f(g(t)) \cdot g'(t)
        \]
    \end{itemize}
\end{frame}

\begin{frame}{Ejemplo en 2D}
    Sea $f(x,y)=x^2+y^2$.
    \[
    \nabla f = (2x,\, 2y)
    \]
    En $(1,1)$: $\nabla f = (2,2)$ (direccion radial).
    \begin{itemize}
        \item Opuesto $(-2,-2)$: direccion de maximo decrecimiento.
    \end{itemize}
\end{frame}

% ============================================================================
% SECCION 3: APLICACION EN IA: DESCENSO DE GRADIENTE
% ============================================================================
\section{Aplicacion en IA: Descenso de Gradiente}

\begin{frame}{Idea fundamental}
    \begin{itemize}
        \item Queremos minimizar la funcion de perdida $\mathcal{L}(\mathbf{w})$.
        \item Como $\nabla\mathcal{L}$ apunta hacia el maximo, para minimizar avanzamos en direccion opuesta:
        \[
        \mathbf{w}_{\text{nuevo}} = \mathbf{w}_{\text{viejo}} - \eta \nabla\mathcal{L}(\mathbf{w}_{\text{viejo}})
        \]
        donde $\eta$ es la \textbf{tasa de aprendizaje}.
    \end{itemize}
\end{frame}

\begin{frame}{Ejemplo: regresion lineal con un parametro}
    Modelo: $\hat{y}=wx$, pérdida MSE para un punto $(x,y)$:
    \[
    \mathcal{L}(w) = (wx - y)^2
    \]
    Gradiente (1D):
    \[
    \frac{d\mathcal{L}}{dw} = 2(wx - y)\,x
    \]
    Actualizacion:
    \[
    w_{\text{nuevo}} = w_{\text{viejo}} - \eta \cdot 2(w_{\text{viejo}}x - y)x
    \]
\end{frame}

\begin{frame}{Ejemplo con dos parametros (intercepto)}
    Modelo: $\hat{y}=w_1x + w_2$, pérdida MSE:
    \[
    \mathcal{L}(w_1,w_2) = (w_1x + w_2 - y)^2
    \]
    Derivadas parciales:
    \[
    \frac{\partial\mathcal{L}}{\partial w_1} = 2(w_1x + w_2 - y)\,x,\quad
    \frac{\partial\mathcal{L}}{\partial w_2} = 2(w_1x + w_2 - y)
    \]
    \begin{align*}
        w_1 &\leftarrow w_1 - \eta\,\frac{\partial\mathcal{L}}{\partial w_1},\\
        w_2 &\leftarrow w_2 - \eta\,\frac{\partial\mathcal{L}}{\partial w_2}
    \end{align*}
\end{frame}

\begin{frame}{Visualizacion del proceso}
    \centering
    \includegraphics[width=0.7\textwidth]{descenso_gradiente.png} % si existe
\end{frame}

% ============================================================================
% SECCION 4: APLICACION EN VISION POR COMPUTADOR: DETECCION DE BORDES
% ============================================================================
\section{Aplicacion en Vision por Computador: Deteccion de Bordes}

\begin{frame}{Gradiente de una imagen}
    \begin{itemize}
        \item Una imagen en grises $I(x,y)$ es una funcion discreta.
        \item Las derivadas parciales se aproximan por diferencias finitas:
        \[
        G_x = I(x+1,y)-I(x,y),\quad G_y = I(x,y+1)-I(x,y)
        \]
        \item Magnitud del gradiente: $G = \sqrt{G_x^2 + G_y^2}$ (alta en bordes).
        \item Direccion: $\theta = \arctan(G_y/G_x)$ (orientacion del borde).
    \end{itemize}
\end{frame}

\begin{frame}{Filtros de Sobel}
    \begin{itemize}
        \item Kernels que aproximan derivadas con suavizado:
        \[
        G_x = \begin{bmatrix}
            -1 & 0 & 1 \\
            -2 & 0 & 2 \\
            -1 & 0 & 1
        \end{bmatrix} * I,\quad
        G_y = \begin{bmatrix}
            -1 & -2 & -1 \\
             0 &  0 &  0 \\
             1 &  2 &  1
        \end{bmatrix} * I
        \]
        \item Mas robustos que diferencias simples.
    \end{itemize}
\end{frame}

% ============================================================================
% SECCION 5: EJEMPLOS CON PYTHON
% ============================================================================
\section{Ejemplos con Python}

\begin{frame}[fragile]{Calculo de gradientes con SymPy}
    \begin{lstlisting}
import sympy as sp

x, y = sp.symbols('x y')
f = x**2 * sp.sin(y) + sp.exp(x*y)

df_dx = sp.diff(f, x)
df_dy = sp.diff(f, y)
print("df/dx =", df_dx)
print("df/dy =", df_dy)

# Evaluar en un punto
valores = {x: 2, y: 1}
print("df/dx(2,1) =", df_dx.subs(valores).evalf())
    \end{lstlisting}
\end{frame}

\begin{frame}[fragile]{Visualizacion de superficie y gradiente}
    \begin{lstlisting}
import numpy as np
import matplotlib.pyplot as plt
from mpl_toolkits.mplot3d import Axes3D

def loss(w1,w2): return (2*w1 + w2 - 5)**2

W1 = np.linspace(-5,5,50)
W2 = np.linspace(-5,5,50)
W1, W2 = np.meshgrid(W1,W2)
Z = loss(W1,W2)

fig = plt.figure(figsize=(12,5))
ax1 = fig.add_subplot(121, projection='3d')
ax1.plot_surface(W1,W2,Z, cmap='viridis', alpha=0.7)
ax1.set_xlabel('w1'); ax1.set_ylabel('w2'); ax1.set_zlabel('Loss')

ax2 = fig.add_subplot(122)
contour = ax2.contour(W1,W2,Z, levels=20, cmap='viridis')
ax2.clabel(contour, inline=True, fontsize=8)
# Dibujar gradiente en algunos puntos
puntos = [(-4,-4), (-2,3), (3,-2), (2,2)]
for w1,w2 in puntos:
    g1 = 4*(2*w1+w2-5)
    g2 = 2*(2*w1+w2-5)
    ax2.arrow(w1,w2, g1*0.1, g2*0.1, head_width=0.3, ec='red')
    ax2.plot(w1,w2,'ro')
ax2.set_xlabel('w1'); ax2.set_ylabel('w2')
plt.tight_layout(); plt.show()
    \end{lstlisting}
\end{frame}

\begin{frame}[fragile]{Descenso de gradiente para regresion lineal}
    \begin{lstlisting}
import numpy as np
import matplotlib.pyplot as plt

np.random.seed(42)
X = 2 * np.random.rand(100, 1)
y = 4 + 3 * X + np.random.randn(100, 1)

X_b = np.c_[np.ones((100, 1)), X]
w = np.random.randn(2, 1)
eta = 0.1
n_iter = 1000
w_history = [w.copy()]

for i in range(n_iter):
    grad = 2/len(X) * X_b.T @ (X_b @ w - y)
    w = w - eta * grad
    w_history.append(w.copy())

print("Intercepto:", w[0][0], "Pendiente:", w[1][0])
    \end{lstlisting}
\end{frame}

\begin{frame}[fragile]{Deteccion de bordes con Sobel}
    \begin{lstlisting}
import numpy as np
import matplotlib.pyplot as plt
from scipy import ndimage
from skimage import data

image = data.camera()

sx = ndimage.sobel(image, axis=0)
sy = ndimage.sobel(image, axis=1)
sobel_mag = np.hypot(sx, sy)

fig, axes = plt.subplots(1,3, figsize=(12,4))
axes[0].imshow(image, cmap='gray')
axes[1].imshow(sobel_mag, cmap='gray')
axes[2].imshow(sx, cmap='gray')
plt.tight_layout(); plt.show()
    \end{lstlisting}
\end{frame}

% ============================================================================
% SECCION 6: NOTACION EN PAPERS CIENTIFICOS
% ============================================================================
\section{Notacion en Papers Cientificos}

\begin{frame}{Notacion comun en IA}
    \begin{itemize}
        \item $\nabla_\theta \mathcal{L}$: gradiente de la perdida respecto a parametros $\theta$.
        \item $\frac{\partial \mathcal{L}}{\partial \mathbf{W}^{(l)}}$: gradiente en la capa $l$.
        \item $\delta^{(l)}$: error propagado hacia atras.
        \item $I_x$, $I_y$: derivadas espaciales de una imagen.
        \item $\nabla^2$: laplaciano (suma de segundas derivadas).
    \end{itemize}
\end{frame}

% ============================================================================
% SECCION 7: EJERCICIOS PROPUESTOS
% ============================================================================
\section{Ejercicios Propuestos}

\begin{frame}{Ejercicios}
    \begin{enumerate}
        \item Calcula las derivadas parciales y el gradiente de $f(x,y,z)=x^2 e^y + \ln(xz) - \sin(yz)$.
        \item Para la funcion de perdida $L(\mathbf{w}) = \frac{1}{2}\|X\mathbf{w}-\mathbf{y}\|^2$, demuestra que $\nabla L = X^T(X\mathbf{w}-\mathbf{y})$.
        \item Implementa descenso de gradiente para la funcion de Himmelblau $f(x,y)=(x^2+y-11)^2+(x+y^2-7)^2$. Visualiza curvas de nivel y trayectoria.
        \item Explica por que el gradiente es perpendicular a las curvas de nivel.
        \item En un paper encuentras $\delta^{(l)} = \nabla_{\mathbf{a}^{(l)}}L \odot \sigma'(\mathbf{z}^{(l)})$. Identifica cada termino.
    \end{enumerate}
\end{frame}

% ============================================================================
% SECCION 8: RESUMEN
% ============================================================================
\section{Resumen}

\begin{frame}{Lo que hemos aprendido}
    \begin{exampleblock}{Conceptos clave}
        \begin{itemize}
            \item \textbf{Derivadas parciales}: sensibilidad individual.
            \item \textbf{Vector gradiente}: direccion de maximo crecimiento, perpendicular a curvas de nivel.
            \item \textbf{Descenso de gradiente}: algoritmo fundamental para minimizar funciones de perdida.
            \item \textbf{Aplicacion en vision}: deteccion de bordes mediante gradiente de imagen.
            \item \textbf{Implementacion}: SymPy para calculo simbolico, NumPy para descenso de gradiente, scikit-image para bordes.
        \end{itemize}
    \end{exampleblock}
\end{frame}

\begin{frame}{Conexion con futuros cursos}
    \begin{itemize}
        \item \textbf{Deep Learning}: backpropagation se basa en gradientes.
        \item \textbf{Optimizacion avanzada}: metodos de segundo orden (Newton, quasi-Newton).
        \item \textbf{Vision por Computador}: flujo optico, descriptores SIFT.
    \end{itemize}
\end{frame}

% --- Diapositiva final ---
\begin{frame}
    \centering
    \Huge \textbf{Gracias!}

\end{frame}

\end{document}